\sectioncentered*{Введение}
\addcontentsline{toc}{section}{Введение}

\subsubsection*{Актуальность темы базы знаний}
    Актуальность разработки базы знаний (БЗ) для интеллектуальной системы подготовки к техническим собеседованиям определяется ростом требований к фундаментальной подготовке IT-специалистов. В условиях конкуренции на рынке труда кандидатам необходимо уверенно владеть базовыми концепциями Computer Science: алгоритмами и структурами данных \cite{Cormen2013}. 

    Для успешного прохождения технических интервью также требуется понимание принципов работы операционных систем и сетевых протоколов \cite{Tanenbaum2012}. 

    При формировании содержания БЗ учитываются практики подготовки и разбора типовых интервью-задач \cite{McDowell2015}. 

    Проект выполняется в рамках дисциплины «Математические основы интеллектуальных систем» и опирается на результаты по инженерии знаний и онтологическому моделированию. Разрабатываемая БЗ поддерживает диагностику уровня знаний, классификацию материалов по доменам и сложности, формирование адаптивных траекторий и генерацию объяснений разной глубины. Целевая аудитория: кандидаты уровней Junior--Senior, специалисты в карьерном переходе и технические менторы \cite{Golenkov2022}. 

\subsubsection*{Цель и задачи курсового проекта}
    Целью курсового проекта является разработка структурированной базы знаний фундаментальных основ Computer Science, выступающей семантическим ядром интеллектуального помощника по подготовке к техническим собеседованиям.

    Для достижения цели решаются следующие задачи:
    \begin{enumerate}
        \item проанализировать предметную область «Фундаментальные основы Computer Science», выделить ключевые концепции, связи и типовые сценарии интервью;
        \item обосновать формализм представления знаний (онтология, метаданные, таксономия тегов) и требования к машиночитаемости, расширяемости и согласованности БЗ;
        \item спроектировать логическую и физическую структуру БЗ с поддержкой навигации, перекрёстных ссылок и основных типов единиц знаний;
        \item реализовать прототип наполнения БЗ материалами по алгоритмам, структурам данных и системным концепциям;
        \item проверить корректность семантических связей, полноту покрытия тем и сопровождаемость структуры.
    \end{enumerate}

\subsubsection*{Объект, предмет и границы разработки}
    \textit{Объектом} разработки является подсистема управления знаниями в составе интеллектуальной системы «Интеллектуальный помощник по подготовке к техническим собеседованиям» (TechInterview), обеспечивающая хранение, организацию и доступ к предметно-ориентированным материалам.

    \textit{Предметом} проектирования выступает база знаний фундаментальных основ Computer Science: виды знаний (декларативные концепции, процедурные паттерны, мета-знания о сложности и связях), формат представления (Markdown с YAML-фронтматтером), логическая модель с доменной таксономией и семантическая навигация.

    \textit{Границы проекта} включают поддержку диагностики знаний, классификацию материалов по темам и уровням, статическую генерацию учебных последовательностей по правилам. Входные данные: структурированные карточки знаний и профили пользователей с указанием целевой роли и уровня. За рамками остаются динамическая адаптация на основе ML-моделей, полноценный пользовательский интерфейс и интеграция с внешними платформами проведения собеседований в реальном времени.

\subsubsection*{Структура пояснительной записки}
    Пояснительная записка состоит из введения, трёх основных разделов и заключения.

    \begin{itemize}
        \item \textit{Раздел 1} содержит анализ предметной области, типовых сценариев, существующих образовательных ресурсов и требований к БЗ.
        
        \item \textit{Раздел 2} описывает проектные решения: онтологический формализм, концептуальную модель, логическую и физическую структуры БЗ, систему метаданных и правила согласованности знаний.
        
        \item \textit{Раздел 3} посвящён реализации: инструментам и форматам, наполнению БЗ типовыми единицами, сценариям использования и результатам тестирования связей.
        
        \item \textit{В заключении} сформулированы выводы о достижении цели проекта, дана оценка качества БЗ (полнота, непротиворечивость, расширяемость) и определены направления дальнейшего развития системы.
    \end{itemize}

    Курсовой проект выполнен самостоятельно, проверен в системе «Антиплагиат». Процент оригинальности соответствует норме, установленной кафедрой ИИТ. Цитирования обозначены ссылками на публикации, указанные в «Списке использованных источников».